\problema{Construct Binary Tree from Preorder and Inorder Traversal}{105}{src/02-arboles/105-construct-binary-tree-from-preorder-and-inorder-traversal.cpp}{M}

Dados dos arreglos de enteros \texttt{preorder} e \texttt{inorder}, donde el primero
es el recorrido en preorden de un árbol binario y el segundo es su recorrido en
inorden, reconstruye y devuelve el árbol binario original. (Se garantiza que no
hay valores duplicados en el árbol).

\paragraph{La idea, con un ejemplo de la vida real.} Un ``recorrido'' es el
orden en que anotamos los nodos cuando paseamos por el árbol. Aquí nos dan dos
listas de ese paseo, y con las dos juntas tenemos que volver a dibujar el árbol,
como armar un mueble con dos hojas de instrucciones distintas. Fíjate en la
forma de cada lista:
\begin{itemize}
    \item \textbf{Preorden:} anota primero la raíz y luego lo demás: \texttt{[RAÍZ, [Subárbol Izquierdo], [Subárbol Derecho]]}. Su primer elemento siempre es el jefe del árbol.
    \item \textbf{Inorden:} anota primero todo lo de la izquierda, luego la raíz y al final lo de la derecha: \texttt{[[Subárbol Izquierdo], RAÍZ, [Subárbol Derecho]]}.
\end{itemize}

Junta las dos pistas y sale magia: \texttt{preorder} te regala quién es la raíz
(su primer elemento). Buscas ese valor en \texttt{inorder} y, como ahí la raíz
está justo en medio, todo lo que quede a su izquierda es el subárbol izquierdo y
todo lo de su derecha es el subárbol derecho. Ya separaste el árbol en dos
mitades más chiquitas.

Por ejemplo, si la raíz es 3 y en \texttt{inorder} aparece en la posición $i$,
entonces hay exactamente $i$ nodos en el subárbol izquierdo. Eso también te dice
que los siguientes $i$ valores de \texttt{preorder} forman esa mitad izquierda, y
el resto la derecha.

\paragraph{Cómo lo resuelve el código, paso a paso.} Repetimos ese mismo truco
una y otra vez sobre pedazos cada vez más pequeños (eso es \emph{recursión}, y
esta estrategia de partir el problema en mitades se llama \emph{dividir y
vencer}):
\begin{enumerate}
    \item Tomamos el siguiente valor de \texttt{preorder}: ese es la raíz del pedazo que estamos armando.
    \item Buscamos dónde está esa raíz en \texttt{inorder} para saber cuántos nodos van a la izquierda y cuántos a la derecha.
    \item Construimos primero el subárbol izquierdo y luego el derecho, cada uno con el mismo procedimiento.
\end{enumerate}
Para que buscar la raíz dentro de \texttt{inorder} sea instantáneo (y no lento,
lo que volvería todo $O(n^2)$), al principio armamos un \emph{mapa hash} (una
agenda que empareja cada valor con su posición) y así encontramos cualquier
índice al instante.

\lstinputlisting{src/02-arboles/105-construct-binary-tree-from-preorder-and-inorder-traversal.cpp}

\complejidad{$O(n)$}{$O(n)$}

\paragraph{¿Qué significa esa complejidad?} $n$ es la cantidad de nodos. El
tiempo $O(n)$ significa que creamos cada nodo una sola vez, y gracias al mapa
hash ya no perdemos tiempo buscando. La memoria $O(n)$ viene de dos cosas: el
propio mapa hash guarda un dato por nodo, y la recursión abre una pila de
llamadas.

\paragraph{Consejos para la entrevista (Amazon).} Es un ejemplo clásico de
``dividir y vencer''. El detalle técnico que más cuenta es pasar el índice de
\texttt{preorder} por \textbf{referencia} (\texttt{int\& preIndex}), para que el
avance de la raíz se mantenga sincronizado entre todas las llamadas recursivas.
Si lo pasas por valor (una copia), los subárboles de la derecha tomarían la raíz
equivocada y el árbol saldría mal.
